% concatenated from DiamondLLocationARXIV050826.tex

\documentclass[11pt,reqno]{amsart}

\usepackage{amsmath,amssymb,amsthm,mathtools}
\usepackage{enumitem}
\usepackage{hyperref}
\usepackage[margin=1in]{geometry}

\newtheorem{theorem}{Theorem}[section]
\newtheorem{proposition}[theorem]{Proposition}
\newtheorem{lemma}[theorem]{Lemma}
\newtheorem{corollary}[theorem]{Corollary}
\theoremstyle{definition}
\newtheorem{definition}[theorem]{Definition}
\newtheorem{example}[theorem]{Example}
\theoremstyle{remark}
\newtheorem{remark}[theorem]{Remark}

\numberwithin{equation}{section}

\title[ Exchangeable Testing Against an Unknown Benchmark]
{Exchangeable Testing Against an Unknown Benchmark}

\author{Alexander Gnedin}
\address{School of Mathematical Sciences, Queen Mary University of London, London E1 4NS, United Kingdom}
\email{a.gnedin@qmul.ac.uk}

\subjclass[2020]{Primary 60J70, 60C05; Secondary 60G09, 62E15}

\keywords{exchangeability, explicit predictive probability, benchmark testing, latent rank, Topp--Leone distribution, P\'olya urn, diffusion limit}

\begin{document}

\begin{abstract}
We generate infinite binary exchangeable sequences by sequential comparison of data points against a latent benchmark. Assuming a prior distribution of the benchmark rank \(R_0\) within an unobserved group, we set up the Bayesian machinery that determines the posterior distribution of the running rank \(R_n\) in purely combinatorial terms. This yields an explicitly computable predictive probability of winning against the benchmark. The normalised running rank converges to a latent strength variable \(X\) with polynomial density, possibly Beta-tilted.

Some min-max tournaments lead to particularly simple multiplicative formulae for predictive probabilities related to priors that generalise the Topp--Leone distribution; for that class we analyse the asymptotics of the associated fixed-\(n\) up-down Markov chains. The limiting diffusion has the classical Wright--Fisher variance but a nonlinear drift expressed explicitly via the prior density of the benchmark.

Mixtures of Beta densities are classical objects in the theory of exchangeable sequences. The contribution of the present work is the combinatorial rank-based updating mechanism and the resulting explicit predictive laws for sequential testing against an unknown benchmark.
\end{abstract}

\maketitle

\section{Introduction}

De Finetti's theorem characterises infinite exchangeable sequences by their mixing measures, thus promising a wide variety of explicit examples already on the level of binary sequences. Despite this, the theory of discrete exchangeable sequences and the hierarchy of processes on combinatorial structures derived from them \cite{Pitman2006,Bertoin2006} is almost entirely built on a single concrete instance, introduced independently by Markov (1917) and Eggenberger--P\'olya (1923) \cite{JohnsonKotz1977}, with roots going back to Bayes and Laplace \cite{Zabell2005}. In its most basic form this is the Markov--P\'olya two-colour urn with the initial composition \((1,1)\), which generalises to \((\alpha,\beta)\)-compositions by conditioning and analytic continuation, and to multi- and infinite-colour processes by the splitting property of the Beta distribution. The exceptional analytic and combinatorial tractability of this class rests on the structure of the binomial moments of the uniform distribution, which, owing to the multiplicative properties of the Beta function, yields an explicit predictive rule. There are many other urn processes exhibiting similar behaviour \cite{FlajoletDumasPuyhaubert2006,Janson2006}; however, they lack the fundamental property of infinite exchangeability, which is known to be equivalent to apparently different forms of probabilistic symmetries \cite{Kallenberg2005,PitmanYor1997,Gnedin2026sym}.

In this paper we depart from the reinforcement paradigm underlying the Markov--P\'olya process, shifting the focus to comparison with latent benchmarks. Such comparisons are the everyday language of stress testing, ranking, quality control, and competitive selection. Structurally, the proposed mathematical model is an infinite sequence in which each element can be uniquely ranked relative to the latent benchmarks and the observed sample data, yielding a uniform distribution for the relative rank, lower ranks corresponding to greater strength. The observer's prior information is encapsulated by the rank distribution within the initial group. Each subsequent data element is then tested against the hidden standards, thereby generating classified outcomes that allow one to update the distribution of the running ranks of the benchmarks. See \cite{Hill2009} for foundational work on rank-based prediction, and \cite{Samuels1991,Gnedin2004} for infinite-arrivals models of sequential choice.

Mixtures of Beta densities (Bernstein polynomials of fixed degree, or their Beta-tilted versions) are classical objects that appear throughout the theory of exchangeable sequences and predictive modelling; see, for instance, the recent review of Fortini \& Petrone \cite{FortiniPetrone2024}. The novelty claimed here is not the class of priors itself, but the combinatorial Bayesian machinery that updates the running rank \(R_n\) of a latent benchmark and the resulting explicit predictive probabilities for sequential testing against that benchmark.

To convey the spirit of the results to follow, and highlight the source that motivated this paper, consider a tournament on a group of four individuals arranged in pairs \((u,v)\), \((w,z)\). The strengths within each pair are compared, winners revealed, and a benchmark \(x\) is chosen as the item weakest among the two pairwise winners. It is readily checked that the rank of \(x\) among \(u,v,w,z\) has distribution \((0,\tfrac23,\tfrac13,0)\), which implies that the next item wins against \(x\) with probability \(\tfrac25\cdot1+\tfrac15\cdot\tfrac13=\tfrac7{15}\) (in which case the running rank of \(x\) increases) or loses with the complementary probability (in which case the running rank of \(x\) does not change). After a series of tests with \(A\) wins and \(B\) failures the predictive probability of success in the next game is
\[
(A,B)\mapsto\frac{(A+2)(A+2B+7)}{(A+B+5)(A+2B+6)}.
\]
This formula is not a quotient of two linear forms characteristic of Markov--P\'olya processes. In accord with de Finetti's law of large numbers the normalised running rank of \(x\) relative to all past data has a limit, which turns out to be a random variable with the Topp--Leone distribution. The latter does not belong to the proper Beta family, but is representable in a standard way as a Beta-mixture, and also as a signed linear combination of Beta densities.

We set up a combinatorial Bayesian framework for analysis of binary testing, classify ranking priors in terms of distinguishability of binary experiments and study fluctuation of the latent running rank \(R_n\). We will show that the birth-death process, identical to the reversible up-down chain on the \(n\)th level of the Pascal graph, is approximable by a Wright--Fisher type diffusion with standard variance but nonlinear drift. The proposed model generalises to more complex exchangeable ranking experiments and ordered partition processes \cite{DongGnedinPitman2007,GnedinPitman2005,GnedinPitman2006}, but here we restrict the primary scope to the basic case of binary sequences.

\section{The binary testing model}

Let an unobserved group of \(k\) past players labelled \(P_{-1},\dots,P_{-k}\) have i.i.d.\ \(\mathrm{U}[0,1]\) latent marks, where a smaller mark signifies a stronger player. We arrange the marks in ranking order and denote them \(V_1<\dots<V_k\). By exchangeability, the correspondence between the ranks of the players relative to this group and the ordered marks is uniformly random. That is, the ranking of \(P_{-1},\dots,P_{-k}\) is independent of the ordered marks. We write
\[
V_r\stackrel{d}{=}\mathrm{Beta}(r,k-r+1)
\]
meaning that \(V_r\) has this Beta density
\[
h_{k,r}(x):=k\binom{k-1}{r-1}x^{r-1}(1-x)^{k-r},\qquad x\in[0,1].
\]

Suppose a choice procedure whose outcome only depends on ranks of the players identifies a winner with rank \(R_0\). The strength of the winner, denoted
\[
X:=V_{R_0},
\]
will be called the benchmark or diamond location. Denote
\[
p_0(r):=\mathbb{P}[R_0=r],\qquad r\in[k],
\]
the distribution of the winner rank relative to \(P_{-1},\dots,P_{-k}\). By independence between the ranks and the marks the pdf of the benchmark is the convex combination of Beta pdf's,
\begin{equation}\label{eq:f0}
f_0(x):=\sum_{r=1}^k p_0(r)\,h_{k,r}(x).
\end{equation}

Now suppose players \(P_1,P_2,\dots\) with i.i.d.\ \(\mathrm{U}[0,1]\) strengths \(U_1,U_2,\dots\) (independent of \(V_1,\dots,V_k\)) sequentially arrive and play against the tournament winner. The outcome of each contest is determined by the rule
\begin{equation}\label{eq:win}
P_n\text{ wins}\quad\Longleftrightarrow\quad U_n<X.
\end{equation}
For \(n=1,2,\dots\) the events \eqref{eq:win} are exchangeable. We assume that the prior information of the observer is the distribution \(p_0\), and that the outcomes \eqref{eq:win} are sequentially observed.

A mental picture is formed by Kingman's paintbox construction \cite{Pitman2006}, where the \(U_i\)'s appear as balls that land either in the box \([0,X]\) to the left of the diamond location or in the box \([X,1]\) to the right, ties having probability zero. The anticipation of the exact diamond location evolves with the count of the balls landing on the left (wins) and on the right (losses).

The key for running the Bayesian prior-to-posterior analysis is the following combinatorial fact, well known in the theory of records and sequential choice, and resulting from the inversion (aka factorial, Lehmer) code of permutations:

\begin{quote}
\textbf{The ranking rule:} the rank of \(P_n\) relative to \(P_{-1},\dots,P_{-k},P_1,\dots,P_n\) has uniform distribution on \(\{1,\dots,n+k\}\); moreover, for \(n=1,2,\dots\) these relative ranks are independent.
\end{quote}

\begin{example}
Let \(w_1,w_2,w_3\) be odds representing a rank prior, the distribution of random variable \(R_0\). The relative rank of \(P_1\) is equally likely to settle in any of four slots, think of inserting \(*\) in the sequence. For instance, the configuration
\[
w_1,*,w_2,w_3
\]
contributes to the winning event \(\tfrac14(w_2+w_3)\). Summing over all possibilities and normalising we see that the player wins with probability \(\mathbb{E}R_0\).
\end{example}

\begin{remark}
The construction of a `uniformly random' ranking of a countable set in terms of the total order on i.i.d.\ marks first appeared in an abstract in the context of the secretary problem \cite{Rubin1966}. Measure-theoretically, the rank rule sets up an isomorphism between the space of sequences of relative ranks of \(P_{-1},\dots,P_{-k},P_1,P_2,\dots\) endowed with the product of discrete uniform distributions and the infinite Lebesgue cube \([0,1]^\infty\) (see \cite{Gnedin2026size} for infinite size-biased models of ranking). Informally, in full analogy with the true probability of heads for an unknown coin, the latent strength of a player appears in the role of the absolute rank in an infinite population.

Similarly to Kingman's paintbox, where balls falling in same box are exchangeable, ranking within the set of winners (respectively, losers) does not carry any predictive utility. Comparison with non-benchmark players does matter, but this does not fit in the framework of binary experiment, and will be considered elsewhere.
\end{remark}

\section{Updating rules}

The benchmark model admits two equivalent descriptions. One, combinatorial, is expressed in terms of the running rank of the benchmark. The other is analytic and is expressed through the latent benchmark location \(X\).

\subsection{The running rank}

Let \(R_n\) denote the rank of the benchmark player relative to
\[
P_{-1},\dots,P_{-k},P_1,\dots,P_n.
\]
After \(A\) wins and \(B=n-A\) losses,
\begin{equation}\label{eq:Rn}
R_n=R_0+A.
\end{equation}
Since \(R_0\in[k]\), the support of \(R_n\) is \(\{A+1,\dots,A+k\}\). The posterior distribution is
\begin{equation}\label{eq:post}
\mathbb{P}[R_n=r+A\mid A,B]=\frac{p_0(r)\,[r]_A^\uparrow\,[k-r+1]_B^\uparrow}{C(A,B)},\qquad r\in[k],
\end{equation}
where the normalising constant is
\begin{equation}\label{eq:C}
C(A,B)=\sum_{j=1}^k p_0(j)\,[j]_A^\uparrow\,[k-j+1]_B^\uparrow.
\end{equation}
The posterior distribution \eqref{eq:post} has a simple combinatorial interpretation. From a state \((r,k-r)\) in the lattice, each win increments the rank coordinate and each loss increments the complementary coordinate. The number of directed paths corresponding to composition \((A,B)\) is the hypergeometric count
\[
\binom{r+A-1}{A}\binom{k-r+B}{B}=\frac{[r]_A^\uparrow\,[k-r+1]_B^\uparrow}{A!\,B!}.
\]
Thus the posterior weight of rank \(r\) is obtained by multiplying the prior weight \(p_0(r)\) by the corresponding path count. Normalisation by the partition function \(C(A,B)\) yields the posterior distribution \eqref{eq:post}.

\subsection{Predictive probability}

When the player \(P_{n+1}\) is introduced, there are \(k+n+1\) possible insertion slots. Conditionally on \(R_n\), exactly \(R_n\) of these slots correspond to a win against the benchmark. Therefore
\begin{equation}\label{eq:piE}
\pi(A,B)=\frac{\mathbb{E}[R_n\mid A,B]}{k+n+1}.
\end{equation}
Since
\[
\mathbb{E}[R_n\mid A,B]=\frac{\sum_{r=1}^k(r+A)p_0(r)\,[r]_A^\uparrow\,[k-r+1]_B^\uparrow}{\sum_{r=1}^k p_0(r)\,[r]_A^\uparrow\,[k-r+1]_B^\uparrow},
\]
we obtain
\begin{equation}\label{eq:pi}
\pi(A,B)=\frac1{k+A+B+1}\frac{\sum_{r=1}^k(r+A)p_0(r)\,[r]_A^\uparrow\,[k-r+1]_B^\uparrow}{\sum_{r=1}^k p_0(r)\,[r]_A^\uparrow\,[k-r+1]_B^\uparrow}.
\end{equation}
Using the elementary identity \((r+A)[r]_A^\uparrow=[r]_{A+1}^\uparrow\), the predictive probability may also be written as
\begin{equation}\label{eq:pi8}
\pi(A,B)=\frac1{k+A+B+1}\frac{\sum_{r=1}^k p_0(r)\,[r]_{A+1}^\uparrow\,[k-r+1]_B^\uparrow}{\sum_{r=1}^k p_0(r)\,[r]_A^\uparrow\,[k-r+1]_B^\uparrow}.
\end{equation}
The simplicity of \eqref{eq:pi8} is one of the main features of the model. The predictor is obtained by shifting a single upper factorial in the numerator.

We now connect to the latent variable. Invoking \eqref{eq:piE},
\begin{equation}\label{eq:piX}
\pi(A,B)=\mathbb{P}(U_{n+1}<X\mid A,B)=\frac{\mathbb{E}[R_n\mid A,B]}{k+n+1}.
\end{equation}
This elementary identity is the key updating rule of the model, hence we formulate it as a theorem.

\begin{theorem}[Rank--diamond location identity]
For every rank prior \(p_0\),
\begin{equation}\label{eq:identity}
\frac{\mathbb{E}[R_n\mid A,B]}{k+n+1}=\mathbb{E}[X\mid A,B].
\end{equation}
Consequently,
\begin{equation}\label{eq:moments}
\pi(A,B)=\mathbb{E}[X\mid A,B]=\frac{\mathbb{E}[X^{A+1}(1-X)^B]}{\mathbb{E}[X^A(1-X)^B]}.
\end{equation}
\end{theorem}

The second assertion goes back to Hausdorff's problem of moments \cite{Hausdorff1921a,Hausdorff1921b}; here it follows from the combinatorial update rule.

\subsection{The posterior distribution of the benchmark}

The posterior pdf of the diamond location \(X\) remains a convex combination of the same Beta distributions,
\begin{equation}\label{eq:fn}
f_n(t\mid A,B)=\sum_{r=1}^k\mathbb{P}[R_n=r+A\mid A,B]\,h_{k,r}(t).
\end{equation}
Comparing with \eqref{eq:f0}, this just shows that the update factors through formula \eqref{eq:post} for the sufficient statistic engaged in a combinatorial random walk on the Pascal graph.

\section{The law of large numbers for running rank}

By the strong version of de Finetti's theorem the proportion of wins in \(n\) exchangeable contests has an a.s.\ limit. Since ranking occurs on the combinatorial level of permutations we may identify this limit random variable with the hypothetical benchmark \(X\). The relative rank remains throughout not exactly observable. Nevertheless, by the virtue of \eqref{eq:Rn} it eventually finds the diamond location:
\begin{equation}\label{eq:LLN}
\frac{R_n}{n}\to X\quad\text{a.s.\ as }n\to\infty.
\end{equation}

\section{Identifiability}

This section is devoted to the analysis of parameters of the binary testing model. The parameters in our model are the size \(k\) of the initial players group and the prior distribution of winner's rank within the group. In this section we denote this discrete prior \((p_1,\dots,p_k)\); the set of such priors is a \(k\)-simplex (of dimension \(k-1\)), denoted \(\mathcal{P}_k\).

It is sometimes convenient to use parallel homogeneous coordinates, passing to odds signature, that is nonzero nonnegative sequences of weights \((w_1,\dots,w_k)\), which normalise by \(W:=\sum_{r=1}^k w_r\). The two pieces of notation are matched as
\[
(w_1,\dots,w_k)\propto(p_1,\dots,p_k),
\]
so we do not distinguish proportional odds signatures.

\begin{definition}
Two rank priors \((p_1,\dots,p_k)\) and \((q_1,\dots,q_\ell)\) are said to be binary-indistinguishable if they induce the (infinite) binary testing procedures with the same distribution. The definition extends to the representation of priors via their odds signatures.
\end{definition}

The partition of priors in indistinguishability classes of this equivalence relation is trivial within simplex \(\mathcal{P}_k\) of a given degree. Indeed, the Beta densities \(h_{k,r}\) of uniform order statistics comprise a positive Bernstein basis of the space of homogeneous polynomials of degree \(k-1\). Thus the set of such mixtures is a \((k-1)\)-dimensional simplex spanned on the extreme points \((h_{k,r},r\in[k])\). By de Finetti's theorem \eqref{eq:LLN}, the binary experiment identifies its continuous prior, hence the binary indistinguishability at fixed level amounts to equality \((p_1,\dots,p_k)=(q_1,\dots,q_k)\).

\begin{definition}
The elements of the simplex \(\mathcal{H}_k\), that is mixtures of Beta distributions \((h_{k,r},r\in[k])\), are called choice priors of degree \(k\).
\end{definition}

The mixing operation for fixed \(k\) is an isomorphism between \(\mathcal{P}_k\) and \(\mathcal{H}_k\), which allowed us on the latent level to distinguish the discrete priors. Now, the simplices of choice priors comprise an inductive system, \(\mathcal{H}_1\subset\mathcal{H}_2\subset\cdots\), where \(\mathcal{H}_1\) has a sole point \(h_{1,1}\) which is the uniform distribution. Looking across degrees \(k\) brings forward a new phenomenon: the discrete uniform priors, with signatures \((1),(1,1),(1,1,1),\dots\), all correspond to the continuous uniform distribution \(\mathrm{U}[0,1]\). Indeed, choosing at random one of the order statistics yields the uniform distribution, hence the basic Markov--P\'olya urn. Thus, geometrically, the barycenters of the simplices \(\mathcal{H}_k\) all coincide with \(h_{1,1}\), while the barycenters of \(\mathcal{P}_k\)'s live in different dimensions despite producing the same binary processes. Probabilistically, this mismatch is easy to explain: testing against a player randomly chosen from \(P_{-1},\dots,P_{-k}\) makes them exchangeable; a phenomenon valid also for finite series of tests, as long as we stay within the paradigm of comparing only with the winner. A more insightful for our study explanation is that the discrete uniform prior \(p_0(r)\) on \(\mathcal{P}_r\) through the ranking rule elevates to the uniform posteriors; in particular if we take a singleton prior group and update the rank of the benchmark through binary comparisons, the observed process will be precisely the same as the basic Markov--P\'olya urn or Kingman's paintbox with uniform split.

Invoking de Finetti's theorem again, the binary indistinguishability of \((w_1,\dots,w_k)\) and \((v_1,\dots,v_\ell)\), \(k<\ell\), amounts to the equality of expansions in the Bernstein bases
\begin{equation}\label{eq:equal}
\sum_{i=1}^k w_i\,h_{k,i}(x)=\sum_{j=1}^\ell v_j\,h_{\ell,j}(x).
\end{equation}
This relation has interpretation in terms of the general Kerov--Vershik theory of multiplicative branching graphs. The instance \cite{VershikKerov1987} in focus here is the Pascal graph canonically associated with the algebra of symmetric functions in two formal variables \(a,b\) via the Hausdorff degree elevation identity \((a+b)a^Ab^B=a^{A+1}b^B+a^Ab^{B+1}\) factored through \(a+b=1\).

Specialising \(a=x\), \(b=(1-x)\) as real numbers leads to the Bernstein basis conversion identity
\begin{equation}\label{eq:Bernstein}
h_{k,i}(x)=\sum_{j=i}^{\ell-k+i}\frac{\binom{j-1}{i-1}\binom{\ell-j}{k-i}}{\binom{\ell}{k}}h_{\ell,j}(x),
\end{equation}
where the hypergeometric coefficients come from the path-counting in the graph. By substituting this identity in the indistinguishability criterion \eqref{eq:equal} on the linearly independent higher-degree basis, we obtain the following combinatorial classification. Two mixtures are identical if and only if the odds signature of the higher degree is the hypergeometric projection of the lower-degree \(w\):
\begin{equation}\label{eq:vj}
v_j=\sum_{i=1}^k w_i\cdot\frac{\binom{j-1}{i-1}\binom{\ell-j}{k-i}}{\binom{\ell}{k}},\qquad\text{for }j=1,\dots,\ell.
\end{equation}

To translate the elevation technique in the language of testing against the benchmark, we introduce the operation of initial group extension.

\begin{definition}
Let \(P_{-k},\dots,P_{-1}\) be a group of players generating some benchmark rank prior \(p_0\). For integer \(\ell>k\), we say that \(P_{-\ell},\dots,P_{-k},\dots,P_{-1}\) is a prior group extension if the players \(P_{-\ell},\dots,P_{-k-1}\) joined the primary group successively, and the distribution of the winner's rank in the whole growing group was updated according to the rank rule. The resulting benchmark rank distribution \(q_0\) on \([\ell]\) is called extension of the prior \(p_0\).
\end{definition}

\begin{theorem}
For \(\ell>k\) two rank priors \(p_0,q_0\) of degree \(k\) and \(\ell\), respectively, are binary-indistinguishable if and only if \(q_0\) is the extension of \(p_0\).
\end{theorem}

\begin{proof}
The Bernstein basis conversion formula shows that the continuous density generated by \(p_0\) of degree \(k\) expands uniquely in the higher-degree basis of degree \(\ell\) with coefficients given by \eqref{eq:vj}. By definition, the extension \(q_0\) is the law of the running rank after the additional \(\ell-k\) players have been inserted according to the ranking rule; the path-counting argument that produces \eqref{eq:vj} is therefore identical to the combinatorial construction of the extension. Hence the two continuous densities coincide if and only if \(q_0\) is that extension.
\end{proof}

\section{Tilting and transforms}

\subsection{Order reversal}

Flipping the interval amounts to measuring strength in the other direction, which transforms a ranking permutation of degree \([k]\) via \(r\mapsto k-r+1\) and makes minimal order statistics maximal. Both orientations turn useful by multiple interval splitting for construction of infinite shuffles \cite{GnedinOlshanski2006,JackaWarren2007}.

\subsection{Beta tilting}

Beta distributions form the family of conjugate priors for homogeneous Bernoulli trials, therefore Beta tilting
\[
f_0(x)\mapsto f_0(x)\,x^{\alpha-1}(1-x)^{\beta-1},\qquad\alpha,\beta\ge0
\]
acts as a shift in the sufficient statistic
\[
(A,B)\mapsto(A+\alpha,B+\beta),
\]
hence all combinatorial formulas remain intact also for fractional tilting variables.

For integer \(A,B\) the tilting is equivalent to passing from the rank prior \(p_0(r)\) to the Bayesian update \eqref{eq:post}. This operation is not equivalent to the prior group extension changing the support of the distribution.

In general, the analytic counterpart of Beta tilting in terms of weights is the Hahn (Beta-Binomial) re-weighting with rising factorials
\[
\tilde{w}_r\propto[r]_{\alpha-1}^\uparrow[k+1-r]_{\beta-1}^\uparrow\,w_r,\qquad r=1,\dots,k.
\]
Setting \(w_1=\dots=w_k\) yields the pure Hahn weights, in which case the predictive probability becomes the classic Markov--P\'olya rule \(\pi(A,B)=(A+1)/(A+B+2)\) regardless of \(k\).

\subsection{The classic process}

The basic \(\mathrm{U}[0,1]\)-driven process corresponds to the weight sequence \(w_1=1\) (i.e.\ a single latent mark that is itself the diamond location); the corresponding homogeneous polynomial of degree \(0\) is the constant \(1\). Therefore the two-parameter Markov--P\'olya process embeds in this picture in different ways.

\begin{enumerate}[label=(\roman*)]
\item by Beta tilting of the uniform prior,
\[
f_0(x)=1\mapsto c\,x^{\alpha-1}(1-x)^{\beta-1},
\]
which yields the classical predictive rule
\[
\pi(A,B)=\frac{A+\alpha}{A+B+\alpha+\beta};
\]
\item by the pure Hahn re-weighting of any finite string of unit weights
\[
(1,\dots,1)\mapsto\bigl([r]_{\alpha-1}^\uparrow[k+1-r]_{\beta-1}^\uparrow\bigr)_{r=1}^k,
\]
which produces the same predictive probabilities (independent of \(k\));
\item A combinatorial tilting amounts to the weight transform of the kind \((1,1,1,1,1,1)\mapsto(0,0,0,1,1,0,0)\), where a contiguous tail block of \(1\)'s is replaced by \(0\)'s. This instance, by the ranking rule, with reduced composition \((0,0,1,0)\), is the \(\mathrm{Beta}(3,2)\) urn.
\end{enumerate}

This has an important consequence related to any dimension of symmetry.

\begin{theorem}
The basic Markov--P\'olya urn is the symmetrisation of an arbitrary weight sequence: for every \(k\), it is the symmetric mixture of the choice model over the orbit of the permutation group of order \(k\) acting on \((w_1,\dots,w_k)\).
\end{theorem}

\begin{proof}
The symmetrisation of an arbitrary sequence of weights is a constant sequence, which yields the uniform density and therefore the classical Markov--P\'olya predictive probabilities.
\end{proof}

From the viewpoint of strength testing, this property characterises the basic Markov--P\'olya process as the test which does not introduce bias among the players.

\section{Examples}

The Markov--P\'olya world is linear-rational. The only known to us example of nonlinearity is a quadratic \cite{Gnedin2010species} appearing as a two-parameter extension of an atom of the Pitman--Yor Chinese Restaurant family. For (untilted) tournament densities the predictive probabilities are rational (quotients of polynomials). This section presents a sample of examples where these polynomials nicely factor. Our champion is the first example of trapezoidal distribution.

\begin{example}
For linear polynomial with rational inhomogeneity
\[
f_0(x)\propto x+p/q
\]
the predictive probability is
\[
\pi(A,B)=\frac{m(A+1,B)}{m(A,B)}=\frac{(A+1)\bigl((q+p)A+pB+2q+3p\bigr)}{(A+B+3)\bigl((q+p)A+pB+q+2p\bigr)}.
\]
\end{example}

\begin{example}
Linear weights \((1,2,3)\):
\begin{equation}\label{eq:123}
\pi(A,B)=\frac{(A+1)(3A+B+7)}{(A+B+3)(3A+B+4)}.
\end{equation}
\end{example}

\begin{example}
Mallows (truncated geometric) weights \((1,q,q^2)\):
\begin{equation}\label{eq:Mallows}
\pi(A,B)=\frac{(A+1)\bigl[q^2A^2+2qAB+(5q^2+2q)A+B^2+(4q+3)B+6q^2+4q+2\bigr]}{(A+B+4)\bigl[q^2A^2+2qAB+(3q^2+2q)A+B^2+(2q+3)B+2q^2+2q+2\bigr]}.
\end{equation}
\end{example}

\begin{example}[Two-stage tournament]
Partition \(k=k_1k_2\) players into \(k_1\) groups of size \(k_2\). In each group take the minimum; among the \(k_1\) group-winners take the maximum. The latent mark \(X\) of the overall winner has CDF and density
\begin{align*}
F(x)&=\bigl(1-(1-x)^{k_2}\bigr)^{k_1},\\
f(x)&=k_1k_2(1-x)^{k_2-1}\bigl(1-(1-x)^{k_2}\bigr)^{k_1-1},\qquad0\le x\le1.
\end{align*}
\end{example}

\begin{example}[Three-stage tournament]
Partition \(k=k_1k_2k_3\) players into \(k_1k_2\) groups of size \(k_3\). In each group take the minimum; partition the resulting winners into \(k_1\) groups of size \(k_2\) and again take the minimum; finally take the maximum among the \(k_1\) remaining candidates. The latent mark \(X\) of the overall winner has CDF
\[
F(x)=\Bigl(1-\bigl(1-(1-x)^{k_3}\bigr)^{k_2}\Bigr)^{k_1},\qquad0\le x\le1.
\]
Such forms are typical for AND/OR trees. The predictive probability factors, but is more cumbersome.
\end{example}

\begin{example}[A tournament tree]
Operations:
\begin{enumerate}[label=(\roman*)]
\item bottom level: two independent max of size \(2\);
\item middle level: one min (bottom-max + singleton) and one max (bottom-max + singleton);
\item root: min of the two middle results.
\end{enumerate}
This is a proper rooted tree with exactly six leaves and three distinct stages. The prior
\[
f(x)\propto x(1-x)^2(1-x^2)(1-x^3)
\]
yields rational product-form
\[
\pi(A,B)=\frac{A+1}{A+B+3}\cdot\frac{A+B+4}{A+2B+6}\cdot\frac{A+2B+7}{A+3B+9}.
\]
\end{example}

\section{Choice operators}

A (single-choice) social choice function is a combinatorial operator which takes a structured finite set (e.g.\ a hypergraph) and outputs its point, a winner. Assuming Plott's path-independence condition this is equivalent to ranking; and the rankings stay consistent as the set grows. This fact motivated the term choice priors.

We isolate next the choice priors of combinatorial nature, which we call pure. Suppose that the winner in the group \(P_{-k},\dots,P_{-1}\) only depends on the ranking permutation; thus a choice function can be represented as a mapping \(C:S_k\to[k]\) assigning to permutation the rank of the winner (not the label of the winner).

\begin{proposition}
Assume that a ranking of the players is uniformly random and that \(R_0\) is the benchmark rank. Then the distribution of \(R_0\) has the form
\begin{equation}\label{eq:pure}
\mathbb{P}[R_0=r]=\frac{w_r}{k!},\qquad r\in[k]
\end{equation}
for some integer composition \((w_1,\dots,w_k)\) of the number \(k!\), and conversely.
\end{proposition}

\begin{proof}
The integer \(w_k\) counts permutations mapped to \(k\).
\end{proof}

We define signatures highlighted by \eqref{eq:pure}, the corresponding rank distributions and elements of \(\mathcal{H}_k\) as pure. Thus there are \(2^{k!}\) pure choice priors in each degree \(k\).

The mapping we just described is the combinatorial ground for the benchmark testing paradigm, explaining the main step away from the Markov--P\'olya world. Permutation viewed as mapping has a natural univariate statistic: the number of cycles. Tilting the number of cycles yields a distinguished one-parameter family of Ewens measures. Further tilting the sizes of cycles adds the Pitman--Yor parameter, and a two-parameter extension by randomisation of the \(\alpha=-1\) case mentioned above.

Permutations viewed as orders carry a single exchangeable measure introduced in \cite{Rubin1966}, though minor steps in the direction of partial exchangeability lead to many parametric families \cite{Gnedin2010coherent,Gnedin2006constrained}; for instance tilting both upper and lower records gives a closest relative of the Ewens measure, however does not allow full symmetry by fixing extreme ranks.

The benchmark processes constructed here is the first step to show that the exchangeable order is nevertheless capable of producing a myriad of fully exchangeable phenomena, being transformed through Plott's social choice operators.

\section{The Topp--Leone family}

We are back to the example mentioned in the Introduction. The basic Topp--Leone density \cite{ToppLeone1955,KotzVandorp2004} with shape parameter \(\nu\) has the form
\begin{equation}\label{eq:TL}
f(x\mid\nu)=2\nu\,x^{\nu-1}(1-x)(2-x)^{\nu-1},\qquad0<x<1,\quad\nu>0,
\end{equation}
having the distribution function
\[
F(x\mid\nu)=x^\nu(2-x)^\nu=\bigl(1-(1-x)^2\bigr)^\nu.
\]
We denote this density \(\mathrm{TL}(\nu)\) and denote \(\mathrm{TL}(\nu,\alpha,\beta)\) its Beta-tilted version, thus \(\mathrm{TL}(\nu,\alpha,\beta)\) has density
\begin{equation}\label{eq:TLtilt}
f(x\mid\nu,\alpha,\beta)\propto f(x\mid\nu)\,x^\alpha(1-x)^\beta,\qquad\alpha>-\nu,\ \beta>-1,
\end{equation}
where the range of admissible parameters is dictated by integrability. See \cite{KotzVandorp2004} for the exciting analysis of modal properties across the range of parameters.

For general \(\nu>0\), \eqref{eq:TL} is not a finite mixture of uniform order-statistic densities of any fixed group size \(k\); computing moments will lead to hypergeometrics. Consequently there is no finite discrete prior \(p_0\) on \([k]\) we need to start with.

However, when \(\nu\) is a positive integer the binomial expansion of \((2-x)^{\nu-1}\) represents \(\mathrm{TL}(\nu)\) as a signed linear combination of Beta kernels, leading to the binomial moments
\begin{equation}\label{eq:mTL}
m(A,B\mid\nu)=2\nu\sum_{i=0}^{\nu-1}\binom{\nu-1}{i}2^{\nu-1-i}(-1)^i\,\mathrm{B}(A+\nu+i,B+2).
\end{equation}
The same expansion is valid for Beta-tilted \(\mathrm{TL}(\nu,\alpha,\beta)\), with \((A,B)\) in the right-hand side replaced by \((A+\alpha,B+\beta)\): \(m(A,B\mid\nu,\alpha,\beta)=m(A+\alpha,B+\beta\mid\nu)\).

The inclusion-exclusion form of \eqref{eq:mTL} is typical for choice priors related to tournament schemes; this one is a \(k=2\nu\) 2-stage tournament of Example 5.

\begin{enumerate}[label=(\roman*)]
\item at the first stage the exchangeable players are organised in pairs, and a weaker player in each pair is revealed,
\item at the second stage, these \(\nu\) weaker players are compared and the tournament winner becomes the strongest of these \(\nu\).
\end{enumerate}

The quotient, via the Beta shift identity \(\mathrm{B}(a+1,c)=\mathrm{B}(a,c)\,a/(a+c)\), simplifies to an explicit rational function of the state. In particular,
\begin{align}
\pi_1(A,B)&=\frac{A+1}{A+B+3},\label{eq:pi1}\\
\pi_2(A,B)&=\frac{(A+2)(A+2B+7)}{(A+B+5)(A+2B+6)},\label{eq:pi2}\\
\pi_3(A,B)&=\frac{(A+3)\bigl(A^2+4AB+4B^2+17A+36B+76\bigr)}{(A+B+7)\bigl(A^2+4AB+4B^2+15A+32B+60\bigr)}.\label{eq:pi3}
\end{align}
The case \(\nu=1\) is the classical \(\mathrm{Beta}(1,2)\) Markov--P\'olya transition rule \eqref{eq:pi1}.

\section{The up-down Markov chain and its diffusion limit}

We now refine the law of large numbers \eqref{eq:LLN} to a diffusion approximation for the fluctuations of the running rank. The construction and the resulting diffusion depend on the prior only through its binomial moment function
\[
m(A,B):=\mathbb{E}[X^A(1-X)^B]=\int_0^1 x^A(1-x)^B f_0(x)\,dx,
\]
already used in \eqref{eq:pi}--\eqref{eq:moments}.

\subsection{The up-down chain}

Fix a sufficiently smooth prior density \(f_0\) on \([0,1]\) and a historic, possibly fractional, composition \((\alpha,\beta)\) with \(m(\alpha,\beta)>0\). Write \(A=\alpha+b\), \(B=\beta+w\) for observed composition \((b,w)\), \(n=b+w\). The marginal law of the number of black balls after \(n\) draws is the mixture of \(\mathrm{Bin}(n,x)\) over \(x\) sampled from \(f_0(\,\cdot\mid\alpha,\beta)\), explicitly
\begin{equation}\label{eq:rho}
\rho_n(b):=\binom{n}{b}\frac{m(A,B)}{m(\alpha,\beta)},\qquad b=0,\dots,n.
\end{equation}
The up-down Markov chain hops between the levels \(n\) and \(n-1\) of the Pascal graph according to the rules:

\begin{enumerate}[label=(\roman*)]
\item \textbf{Down move:} a ball is chosen uniformly at random among the \(n\) balls and removed, leaving either \((b-1,w)\) or \((b,w-1)\). This is the standard backward transition.
\item \textbf{Up move:} a black or white ball is added according to the particular transition law \(\pi(A,B)=m(A+1,B)/m(A,B)\) of \eqref{eq:moments}.
\end{enumerate}

The combination of the two moves preserves \(n=b+w\) and defines a birth--death chain on the level set \(\{(b,w):b+w=n\}\). The transition probabilities out of the composition \((b,w)\) are
\begin{align}
p_+(b,w)&:=\mathbb{P}[(b,w)\to(b+1,w-1)]=\frac{w}{n}\pi(b,w-1),\label{eq:p+}\\
p_-(b,w)&:=\mathbb{P}[(b,w)\to(b-1,w+1)]=\frac{b}{n}\bigl(1-\pi(b-1,w)\bigr),\label{eq:p-}\\
p_0(b,w)&:=\mathbb{P}[(b,w)\to(b,w)]=1-p_+(b,w)-p_-(b,w),\label{eq:p0}
\end{align}
interpreted as in the case of the removed ball's colour and its replacement, as before.

\begin{remark}
When \(f_0\) is the mixture \eqref{eq:f0} for a finite degree rank prior \(p_0\) on \([k]\), the up-move is literally the combinatorial ranking-rule update of Section 2: \(\pi(b,w-1)\) equals the predictive probability \eqref{eq:pi} at the reduced composition, and \(m(A,B)\) is, up to normalisation, the partition function \(C(A,B)\) of \eqref{eq:C}. This is the combinatorial content underlying the up-down chain in general. The Topp--Leone family is the case where, for non-integer \(\nu\), \(f_0\) is not of this finite-mixture form, so the chain must be described directly through \(m(A,B)\); for integer \(\nu\) it reverts to the finite combinatorial picture.
\end{remark}

The next fact is an instance of a very general result. We give a complete proof to review some historic perspective.

\begin{lemma}[Reversibility]
For every prior density \(f_0\) and every admissible historic composition \((\alpha,\beta)\), the up-down chain on \(\{(b,w):b+w=n\}\) is reversible with respect to the law \(\rho_n\) of \eqref{eq:rho}:
\[
\rho_n(b)\,p_+(b,w)=\rho_n(b+1)\,p_-(b+1,w-1),\qquad b=0,\dots,n-1.
\]
\end{lemma}

\begin{proof}
Write \(A=\alpha+b\), so that with \(B=\beta+w\),
\[
\pi:=\pi(b,w-1)=\frac{m(A+1,B-1)}{m(A,B-1)}
\]
is the value appearing in both transition probabilities,
\[
p_+(b,w)=\frac{w}{n}\pi,\qquad p_-(b+1,w-1)=\frac{b+1}{n}(1-\pi).
\]
By \eqref{eq:rho}, \(\rho_n(b)=\binom{n}{b}m(A,B)/m(\alpha,\beta)\) and \(\rho_n(b+1)=\binom{n}{b+1}m(A+1,B-1)/m(\alpha,\beta)\). The Hausdorff \cite{Hausdorff1921a,Hausdorff1921b} recursion for binomial moments, valid for any probability distribution on \([0,1]\),
\begin{equation}\label{eq:Hausdorff}
m(A,B-1)=m(A,B)+m(A+1,B-1),
\end{equation}
gives directly \(1-\pi=m(A,B)/m(A,B-1)\). Together with
\[
\binom{n}{b}\frac{w}{n}=\binom{n}{b+1}\frac{b+1}{n}=\binom{n-1}{b},
\]
this yields
\[
\rho_n(b)\,p_+(b,w)=\binom{n-1}{b}(1-\pi)\frac{m(A+1,B-1)}{m(\alpha,\beta)}=\rho_n(b+1)\,p_-(b+1,w-1).
\]
\end{proof}

\begin{remark}
Reversibility is common for up-down chains of Diaconis and Fill \cite{DiaconisFill1990} and their combinatorial relatives on branching graphs \cite{BorodinOlshanski2009,DongGnedinPitman2007,GnedinPitman2006,Petrov2013,Petrov2009}. Nothing in the proof uses a specific \(f_0\): only the Hausdorff recursion \eqref{eq:Hausdorff}, which holds for every probability distribution on \([0,1]\).
\end{remark}

\subsection{Diffusion limit}

To scale the up-down chain, represented as the birth--death process on the \(n\)th level \(\{(b,w):b+w=n\}\), set \(x=b/n\in[0,1]\) and accelerate time by the factor \(n^2\). As \(n\to\infty\) with \(b/n\to x\), the arguments of \(\pi\) in \eqref{eq:p+}--\eqref{eq:p-} both diverge with ratio tending to \(x\). Since \(\pi(A,B)=\mathbb{E}[X\mid A,B]\) is the posterior mean of \(X\) after observing a fraction \(x\) of black balls in a large number of trials, Bayesian consistency -- a form of de Finetti's theorem, valid whenever \(f_0\) is continuous and positive at \(x\) -- gives
\begin{equation}\label{eq:piinf}
\pi^\infty(x):=\lim_{n\to\infty}\pi(\lfloor xn\rfloor,\lfloor(1-x)n\rfloor)=x,\qquad0<x<1,
\end{equation}
for every admissible prior \(f_0\). Because the leading terms of \(p_+(b,w)\) and \(p_-(b,w)\) both converge to \(x(1-x)\), they cancel at order \(1\), and the nontrivial drift emerges only at the next order, \(O(1/n)\). Similarly,
\begin{equation}\label{eq:sigma}
\sigma^2(x):=\lim_{n\to\infty}n\bigl(p_+(b,w)+p_-(b,w)\bigr)=2x(1-x),
\end{equation}
independently of \(f_0\), because to leading order both terms in \eqref{eq:p+}--\eqref{eq:p-} tend to \(x(1-x)\): the diffusion coefficient is universal, and only the drift will depend on \(f_0\).

Rather than expand \(\pi\) to order \(1/n\) directly, it is more spare to determine the limiting drift from Lemma~1 together with the classical correspondence between reversible one-dimensional diffusions and their speed densities (cf.\ \cite{EthierKurtz1986}).

\begin{proposition}
For \(x=\lfloor xn\rfloor/n\) fixed, \(n\rho_n(\lfloor xn\rfloor)\to f_0(x)\) as \(n\to\infty\).
\end{proposition}

\begin{proof}
By \eqref{eq:rho}, \(\rho_n(b)=\binom{n}{b}\int_0^1 x^b(1-x)^{n-b}f_0(x\mid\alpha,\beta)\,dx\) is the \(b\)th weight of the Bayes mixture of \(\mathrm{Bin}(n,x)\) over \(x\sim f_0(\,\cdot\mid\alpha,\beta)\), and is thus governed by Hausdorff's moment problem \cite{Hausdorff1921a,Hausdorff1921b} for this mixing measure. By the local central limit theorem, the binomial kernel \(\binom{n}{b}x^b(1-x)^{n-b}\), as a function of \(x\), concentrates on a window of width \(O(n^{-1/2})\) about \(x=b/n\), with total mass \(1\); consequently \(n\binom{n}{b}x^b(1-x)^{n-b}\,dx\to\delta_x(dx)\) weakly, and integrating the continuous density \(f_0(\,\cdot\mid\alpha,\beta)\) against this approximate identity gives \(n\rho_n(b)\to f_0(x\mid\alpha,\beta)\).
\end{proof}

Combined with Lemma~1, Proposition~2 identifies the stationary/reversible density of the diffusion limit: it must be \(f_0(\,\cdot\mid\alpha,\beta)\). For a generator \(\mathcal{L}f=a(x)f''+\mu(x)f'\) reversible with respect to a density \(\pi\), the stationarity equation \((a\pi)'=\mu\pi\) (obtained by a single integration of the forward equation \((a\pi)''-(\mu\pi)'=0\), with natural boundary conditions on \([0,1]\)) determines the drift once \(a\) and \(\pi\) are known:
\begin{equation}\label{eq:mu}
\mu(x)=a'(x)+a(x)\bigl(\log\pi(x)\bigr)'.
\end{equation}
Here \(a(x)=x(1-x)\) by \eqref{eq:sigma} (writing the generator as \(\mathcal{L}f=x(1-x)f''+\mu f'\), i.e.\ \(\sigma^2=2a\)), and \(\pi(x)=f_0(x\mid\alpha,\beta)\). Substituting into \eqref{eq:mu} yields the following general theorem.

\begin{theorem}[Diffusion limit]
Let \(f_0\) be a prior density on \([0,1]\), differentiable and positive on \((0,1)\) with \((\log f_0)'\) locally integrable there. Started from a composition with historic parameters \((\alpha,\beta)\), the rescaled up-down chain
\[
Y_n(t):=n^{-1}B_{\lfloor n^2 t\rfloor}
\]
converges weakly, as \(n\to\infty\), to the \([0,1]\)-valued diffusion with generator
\begin{equation}\label{eq:gen}
\mathcal{L}_{f_0,\alpha,\beta}f(x)=x(1-x)f''(x)+\mu_{f_0}(x;\alpha,\beta)f'(x),
\end{equation}
where
\begin{equation}\label{eq:mu0}
\mu_{f_0}(x;\alpha,\beta)=(1-2x)+x(1-x)\bigl(\log f_0(x\mid\alpha,\beta)\bigr)'.
\end{equation}
\end{theorem}

Equivalently, \(Y\) solves the stochastic differential equation
\begin{equation}\label{eq:SDE}
\mathrm{d}Y_t=\mu_{f_0}(Y_t;\alpha,\beta)\,\mathrm{d}t+\sqrt{2Y_t(1-Y_t)}\,\mathrm{d}B_t.
\end{equation}
The process \(Y\) is reversible with respect to \(f_0(\,\cdot\mid\alpha,\beta)\), which is its unique stationary density on \((0,1)\). The diffusion coefficient \(2x(1-x)\) is universal, the same Wright--Fisher form for every admissible \(f_0\); only the drift \eqref{eq:mu0} depends on the prior, and it does so only through \((\log f_0)'\).

This is the general form: every up-down chain built from a smooth prior \(f_0\) has a Wright--Fisher-type diffusion limit, with the same \(2x(1-x)\) variance and a drift that is completely explicit once \(f_0\) is known. The Topp--Leone family is a convenient worked example -- convenient because \(f_0=f_\nu(\,\cdot\mid\alpha,\beta)\) has an explicit elementary form and (for integer \(\nu\)) a rational pre-limit transition rule -- but Theorem~4 applies verbatim to any other smooth prior, e.g.\ the finite tournament priors of Sections~2 and the min-max constructions once extended off the lattice by \eqref{eq:f0}, or the linear and geometric-weight examples of Section~7.

\section{Diffusion limit for the Topp--Leone family}

Theorem~4 applies to the up-down chain simply by substituting \(f_0=f_\nu(\,\cdot\mid\alpha,\beta)\propto x^{\nu+\alpha-1}(1-x)^\beta(2-x)^{\nu-1}\): nothing further needs to be proved. Since
\[
\bigl(\log f_\nu(x\mid\alpha,\beta)\bigr)'=\frac{\nu+\alpha-1}{x}-\frac{\beta}{1-x}-\frac{\nu-1}{2-x},
\]
formula \eqref{eq:mu0} gives, after simplification, the drift term
\begin{equation}\label{eq:munu}
\mu_\nu(x;\alpha,\beta)=(\alpha+\nu)-(\alpha+\beta+\nu+2)x-(\nu-1)\frac{x(1-x)}{2-x}.
\end{equation}
Equivalently, writing \(Y\) for the diffusion limit of Theorem~4 with this \(\mu_\nu\), \(Y\) solves the stochastic differential equation
\begin{equation}\label{eq:SDEnu}
\mathrm{d}Y_t=\mu_\nu(Y_t;\alpha,\beta)\,\mathrm{d}t+\sqrt{2Y_t(1-Y_t)}\,\mathrm{d}B_t,
\end{equation}
reversible with respect to \(f_\nu(\,\cdot\mid\alpha,\beta)\), its unique stationary density on \((0,1)\). What is special to integer \(\nu\) is only that the pre-limit transition rule \(\pi_\nu(A,B)\) and the invariant law \(\rho_n\) of \eqref{eq:rho} are then given by the explicit rational and hypergeometric formulas; the diffusion limit itself, and formula \eqref{eq:munu}, hold for every real \(\nu>0\).

\begin{example}[Explicit cases]
For \(\nu=1,2,3\), \eqref{eq:munu} gives
\begin{align*}
\mu_1(x;\alpha,\beta)&=(\alpha+1)-(\alpha+\beta+3)x,\\
\mu_2(x;\alpha,\beta)&=(\alpha+2)-(\alpha+\beta+4)x-\frac{x(1-x)}{2-x},\\
\mu_3(x;\alpha,\beta)&=(\alpha+3)-(\alpha+\beta+5)x-2\frac{x(1-x)}{2-x}.
\end{align*}
The formula for \(\mu_1\) reproduces the drift of the Wright--Fisher diffusion with two-sided immigration at rates \(\alpha+1,\beta+2\). The base Topp--Leone mixing corresponds to the initial parameters \(\alpha=\beta=0\).

For \(\nu\ge2\) the additional term
\[
-(\nu-1)\frac{x(1-x)}{2-x}
\]
is a bounded perturbation of the linear immigration drift, vanishing at both endpoints of \([0,1]\), so the boundaries behave as in the Wright--Fisher case. Because the diffusion coefficient \(x(1-x)\) itself vanishes at \(0\) and \(1\), the sign of \(\mu_\nu(x;\alpha,\beta)\) alone does not indicate the direction in which stationary mass accumulates: by \eqref{eq:mu}, the stationary log-density satisfies
\[
\bigl(\log f_\nu(x\mid\alpha,\beta)\bigr)'=\frac{\mu_\nu(x;\alpha,\beta)-(1-2x)}{x(1-x)},
\]
so that \(f_\nu(\,\cdot\mid\alpha,\beta)\) increases where \(\mu_\nu(x;\alpha,\beta)>1-2x\) and decreases where \(\mu_\nu(x;\alpha,\beta)<1-2x\): the relevant comparison is with \(1-2x=a'(x)\), not with \(0\). The unique mode \(p^*\) of the following lemma is exactly the root in \((0,1)\) of \(\mu_\nu(x;\alpha,\beta)=1-2x\), when one exists.
\end{example}

\begin{lemma}
For \(\nu>0\) and all admissible \(\alpha>-\nu\), \(\beta>-1\), the density \(f_\nu(\,\cdot\mid\alpha,\beta)\) has a unique mode \(p^*\in[0,1)\): it is interior if \(\nu+\alpha>1\), and equals \(0\) (with \(f_\nu(\,\cdot\mid\alpha,\beta)\) strictly decreasing on \((0,1)\)) if \(\nu+\alpha\le1\); in the latter case the density is moreover unbounded as \(p\to0\) when \(\nu+\alpha<1\). In every case \(p^*<1\): since \(\beta+1>0\), mass can accumulate only at the left endpoint, never at the right.
\end{lemma}

\begin{proof}
Differentiating the loglikelihood
\[
\log f_\nu(p\mid\alpha,\beta)=\mathrm{const}+(\nu+\alpha-1)\log p+(\beta+1)\log(1-p)+(\nu-1)\log(2-p)
\]
and clearing denominators shows that the stationary points of \(p\mapsto f_\nu(p\mid\alpha,\beta)\) in \((0,1)\) are the roots of the quadratic \(\varphi(p)=ap^2+bp+c\), where
\begin{align*}
a&=2\nu+\alpha+\beta-1,\\
b&=-4\nu-3\alpha-2\beta+2,\\
c&=2\nu+2\alpha-2.
\end{align*}
As \(p\to0+\), \(f_\nu(p\mid\alpha,\beta)\sim\mathrm{const}\cdot p^{\nu+\alpha-1}\), which vanishes, tends to a finite positive limit, or diverges according as \(\nu+\alpha-1\) is positive, zero, or negative (the singular case remaining integrable since \(\alpha>-\nu\)); as \(p\to1-\), \(f_\nu(p\mid\alpha,\beta)\sim\mathrm{const}\cdot(1-p)^{\beta+1}\to0\), because \(\beta+1>0\) throughout the admissible range. Direct examination of \(\varphi\) on \((0,1)\) then shows it has exactly one root there when \(\nu+\alpha>1\), giving the interior mode, and none when \(\nu+\alpha\le1\), in which case \(f_\nu(\,\cdot\mid\alpha,\beta)\) is strictly decreasing throughout and \(p^*=0\).
\end{proof}

See \cite{KotzVandorp2004} for other parameter regimes.

\begin{example}
For \(\nu=1\), \(\alpha=\beta=0\) (base \(\mathrm{TL}(1)=\mathrm{Beta}(1,2)\)), the drift \(\mu_1(x)=1-3x\) is positive for \(x<1/3\): read naively, it would suggest mass being pushed away from \(0\) on that range. Yet the mode is at \(p^*=0\) by Lemma~2 (here \(\nu+\alpha=1\)), because the correct comparison gives \(\mu_1(x)-(1-2x)=-x<0\) throughout \((0,1)\): the stationary density \(f_1(x\mid0,0)=2(1-x)\) is strictly decreasing everywhere, consistently, even though \(\mu_1(x)\) itself changes sign at \(x=1/3\). The same phenomenon persists for every \(\nu\ge1\) whenever \(\nu+\alpha\le1\): the density's mass is pushed to the left towards \(p^*=0\) not because the raw drift is negative there, but because the vanishing diffusion coefficient at the boundary makes \(\mu_\nu(x)-(1-2x)\), rather than \(\mu_\nu(x)\) itself, the operative quantity.
\end{example}

\begin{remark}
Theorem~4, applied via \eqref{eq:mu0}, makes no reference to \(\nu\) being an integer: formula \eqref{eq:munu} is valid for every real \(\nu>0\), even though the finite-\(n\) transition rule \(\pi_\nu(A,B)\) is then genuinely hypergeometric rather than rational. The restriction to integer \(\nu\) is needed only for the closed rational form of the pre-limit chain.
\end{remark}

\begin{thebibliography}{99}

\bibitem{AlsmeyerIksanovMarynych2017}
G.~Alsmeyer, A.~Iksanov and A.~Marynych,
Functional limit theorems for the number of occupied boxes in the Bernoulli sieve,
\textit{Stochastic Processes and their Applications} \textbf{127} (2017) 985--1007.

\bibitem{BraganetsIksanov2025}
O.~Braganets and A.~Iksanov,
On intermediate levels of nested occupancy scheme in random environment generated by stick-breaking: the case of heavy tails,
\textit{Probability and Mathematical Statistics} \textbf{45} (2025) 1--25.

\bibitem{JohnsonKotz1977}
N.~L.~Johnson and S.~Kotz,
\textit{Urn Models and Their Application},
Wiley, New York, 1977.

\bibitem{Zabell2005}
S.~L.~Zabell,
\textit{Symmetry and its Discontents: Essays on the History of Inductive Probability},
Cambridge University Press, New York, 2005.

\bibitem{KotzVandorp2004}
S.~Kotz and J.~R.~van Dorp,
\textit{Beyond Beta: Other Continuous Families of Distributions with Bounded Support and Applications},
World Scientific, Singapore, 2004.

\bibitem{ToppLeone1955}
C.~W.~Topp and F.~C.~Leone,
A family of J-shaped frequency functions,
\textit{Journal of the American Statistical Association} \textbf{50} (1955) 209--219.

\bibitem{DiaconisFill1990}
P.~Diaconis and J.~A.~Fill,
Strong stationary times via a new form of duality,
\textit{Annals of Probability} \textbf{18} (1990) 1483--1522.

\bibitem{EthierKurtz1986}
S.~N.~Ethier and T.~G.~Kurtz,
\textit{Markov Processes: Characterization and Convergence},
Wiley, New York, 1986.

\bibitem{BorodinOlshanski2009}
A.~Borodin and G.~Olshanski,
Infinite-dimensional diffusions as limits of random walks on partitions,
\textit{Probability Theory and Related Fields} \textbf{144} (2009) 281--318.

\bibitem{DongGnedinPitman2007}
R.~Dong, A.~Gnedin and J.~Pitman,
Exchangeable partitions derived from Markovian coalescents,
\textit{Annals of Applied Probability} \textbf{17} (2007) 1172--1201.

\bibitem{FlajoletDumasPuyhaubert2006}
P.~Flajolet, P.~Dumas and V.~Puyhaubert,
Some exactly solvable models of urn process theory,
\textit{Discrete Mathematics \& Theoretical Computer Science} \textbf{AG} (2006) 59--118.

\bibitem{GnedinIksanovMarynych2010}
A.~Gnedin, A.~Iksanov and A.~Marynych,
The Bernoulli sieve: an overview,
\textit{Discrete Mathematics \& Theoretical Computer Science}, DMTCS Proceedings vol.~AM, 21st International Meeting on Probabilistic, Combinatorial, and Asymptotic Methods in the Analysis of Algorithms (AofA'10) (2010) 429--442.

\bibitem{GnedinOlshanski2006}
A.~Gnedin and G.~Olshanski,
Coherent permutations with descent statistic and the boundary problem for the graph of zigzag diagrams,
\textit{International Mathematics Research Notices} \textbf{2006} (2006), Article ID 51968, 39~pp.

\bibitem{JackaWarren2007}
S.~Jacka and J.~Warren,
Random orderings of the integers and card shuffling,
\textit{Stochastic Processes and their Applications} \textbf{117} (2007) 708--719.

\bibitem{GnedinPitman2006}
A.~Gnedin and J.~Pitman,
Exchangeable Gibbs partitions and Stirling triangles,
\textit{Journal of Mathematical Sciences} \textbf{138} (2006) 5674--5685.

\bibitem{Janson2006}
S.~Janson,
Limit theorems for triangular urn schemes,
\textit{Probability Theory and Related Fields} \textbf{134} (2006) 417--452.

\bibitem{Gnedin2010species}
A.~Gnedin,
A species sampling model with finitely many types,
\textit{Electronic Communications in Probability} \textbf{15} (2010) 79--88.

\bibitem{GnedinPitman2005}
A.~Gnedin and J.~Pitman,
Regenerative composition structures,
\textit{Annals of Probability} \textbf{33} (2005) 445--479.

\bibitem{Hausdorff1921a}
F.~Hausdorff,
Summationsmethoden und Momentfolgen.~I,
\textit{Mathematische Zeitschrift} \textbf{9} (1921) 74--109.

\bibitem{Hausdorff1921b}
F.~Hausdorff,
Summationsmethoden und Momentfolgen.~II,
\textit{Mathematische Zeitschrift} \textbf{9} (1921) 280--299.

\bibitem{Gnedin2010coherent}
A.~Gnedin,
Coherent random permutations with biased record statistics,
\textit{Discrete Mathematics} \textbf{310} (2010) 3114--3123.

\bibitem{Gnedin2006constrained}
A.~Gnedin,
Constrained exchangeable partitions,
in \textit{Proceedings of the Fourth Colloquium on Mathematics and Computer Science}, DMTCS Proceedings vol.~AG (2006) 391--400.

\bibitem{Iksanov2022}
A.~Iksanov,
On intermediate levels of a nested occupancy scheme in a random environment generated by stick-breaking,
\textit{Journal of Theoretical Probability} \textbf{35} (2022) 2145--2178.

\bibitem{Pitman2006}
J.~Pitman,
\textit{Combinatorial Stochastic Processes},
Lecture Notes in Mathematics, Vol.~1875, Springer-Verlag, Berlin, 2006.

\bibitem{Petrov2013}
L.~Petrov,
\(sl(2)\) operators and Markov processes on branching graphs,
\textit{Journal of Algebraic Combinatorics} \textbf{38} (2013) 553--609.

\bibitem{Petrov2009}
L.~Petrov,
Two-parameter family of infinite-dimensional diffusions on the Kingman simplex,
\textit{Functional Analysis and Its Applications} \textbf{43} (2009) 279--296.

\bibitem{Gnedin2026size}
A.~Gnedin,
Infinite size-biased orders,
preprint, 2026.

\bibitem{Hill2009}
T.~P.~Hill,
Knowing when to stop,
\textit{American Scientist} \textbf{97} (2009) 126--133.

\bibitem{Samuels1991}
S.~M.~Samuels,
Secretary problems,
in \textit{Handbook of Sequential Analysis}, B.~K.~Ghosh and P.~K.~Sen (eds), Marcel Dekker, 1991, 381--405.

\bibitem{Gnedin2004}
A.~Gnedin,
Best choice from the planar Poisson process,
\textit{Stochastic Processes and their Applications} \textbf{111} (2004) 233--254.

\bibitem{VershikKerov1987}
A.~M.~Vershik and S.~V.~Kerov,
Locally semisimple algebras. Combinatorial theory and the \(K_0\)-functor,
\textit{Journal of Soviet Mathematics} \textbf{38} (1987) 1701--1733.

\bibitem{Bertoin2006}
J.~Bertoin,
\textit{Random Fragmentation and Coagulation Processes},
Cambridge University Press, 2006.

\bibitem{Kallenberg2005}
O.~Kallenberg,
\textit{Probabilistic Symmetries and Invariance Principles},
Springer, 2005.

\bibitem{PitmanYor1997}
J.~Pitman and M.~Yor,
The two-parameter Poisson--Dirichlet distribution derived from a stable subordinator,
\textit{Annals of Probability} \textbf{25} (1997) 855--900.

\bibitem{Gnedin2026sym}
A.~Gnedin,
Symmetries of random partitions,
preprint, 2026.

\bibitem{Rubin1966}
H.~Rubin,
The `secretary' problem,
\textit{Annals of Mathematical Statistics} \textbf{37} (1966) 544 (abstract).

\bibitem{FortiniPetrone2024}
S.~Fortini and S.~Petrone,
Exchangeability, prediction and predictive modeling in Bayesian statistics,
arXiv:2402.10126v2, 2024.

\end{thebibliography}

\end{document}
```
